% Reviewed translation: PDF pages 139--146 / printed pages 125--132.
% Ends before Section 2 begins on the shared PDF page 146.

\setcounter{statement}{7}
\begin{Theorem}\label{thm:primitive-homogeneous-stokes-convergence}
在 Hypotheses H1, H2 和 H3 下, 问题
\eqref{eq:primitive-stokes-discrete-mixed} 有唯一解
$(\mathbf u_h,p_h)\in V_h\times M_h$, 且 $\mathbf u_h$ 也是问题
\eqref{eq:primitive-stokes-discrete-kernel-problem} 的唯一解. 此外,
$(\mathbf u_h,p_h)$ 收敛到问题
\eqref{eq:primitive-homogeneous-stokes-strong} 的解 $(\mathbf u,p)$:
\begin{equation}\label{eq:primitive-stokes-convergence}
\lim_{h\to0}\left\{
|\mathbf u_h-\mathbf u|_{1,\Omega}
+\|p_h-p\|_{0,\Omega}
\right\}=0.
\end{equation}
进一步, 若对某个满足 $1\leq m\leq l$ 的整数 $m$ 有
\[
(\mathbf u,p)\in H^{m+1}(\Omega)^N
\times(H^m(\Omega)\cap L_0^2(\Omega)),
\]
则有误差界
\begin{equation}\label{eq:primitive-stokes-h1-pressure-error}
|\mathbf u-\mathbf u_h|_{1,\Omega}
+\|p-p_h\|_{0,\Omega}
\leq Ch^m\left\{
\|\mathbf u\|_{m+1,\Omega}+\|p\|_{m,\Omega}
\right\}.
\end{equation}
\end{Theorem}

\begin{Proof}
应用 Theorem~\ref{thm:primitive-abstract-discrete-error}.
由 Hypothesis H3, 空间对 $(X_h,M_h)$ 满足一致 inf-sup 条件,
所以只需检验 $a(\cdot,\cdot)$ 的椭圆性, 即可得到问题
\eqref{eq:primitive-stokes-discrete-mixed} 有唯一解.
当 $a(\cdot,\cdot)$ 由 \eqref{eq:primitive-stokes-bilinear-forms} 的 (b) 定义时,
\[
a(\mathbf v,\mathbf v)=\nu|\mathbf v|_{1,\Omega}^2
\qquad\forall\mathbf v\in H^1(\Omega)^N.
\]
当 $a(\cdot,\cdot)$ 由 \eqref{eq:primitive-stokes-bilinear-forms} 的 (a) 定义时,
使用 Korn 不等式, 得到
\[
a(\mathbf v,\mathbf v)\geq\nu|\mathbf v|_{1,\Omega}^2
\qquad\forall\mathbf v\in H_0^1(\Omega)^N.
\]

\MathBookSourcePage{140}
因此在两种情形中, 形式 $a(\cdot,\cdot)$ 都在 $X$ 上椭圆.
Theorem~\ref{thm:primitive-abstract-discrete-error} 于是表明问题
\eqref{eq:primitive-stokes-discrete-mixed} 有唯一解
$(\mathbf u_h,p_h)\in X_h\times M_h$, 其中 $\mathbf u_h$ 是问题
\eqref{eq:primitive-stokes-discrete-kernel-problem} 的唯一解, 且
\begin{equation}\label{eq:primitive-stokes-best-approximation}
\begin{split}
|\mathbf u-\mathbf u_h|_{1,\Omega}+\|p-p_h\|_{0,\Omega}
\leq C_1\biggl\{&
\inf_{\mathbf v_h\in X_h}|\mathbf u-\mathbf v_h|_{1,\Omega}\\
&+\inf_{q_h\in M_h}\|p-q_h\|_{0,\Omega}
\biggr\},
\end{split}
\end{equation}
其中常数 $C_1$ 与 $h$ 无关.

注意, 若 $S_h$ 不把 $L_0^2(\Omega)$ 映到 $M_h$ 上, 可以用
\[
\bar S_hq=S_hq-\frac{1}{\operatorname{meas}(\Omega)}
\int_\Omega S_hq\,dx
\]
代替它. 此时 $\bar S_h\in\mathcal L(L^2(\Omega);M_h)$, 并且
\[
\|\bar S_hq-q\|_{0,\Omega}\leq\|S_hq-q\|_{0,\Omega}
\qquad\forall q\in L_0^2(\Omega).
\]
因此, 若 $p\in H^m(\Omega)\cap L_0^2(\Omega)$, Hypothesis H2 给出
\[
\inf_{q_h\in M_h}\|p-q_h\|_{0,\Omega}
\leq\|p-S_hp\|_{0,\Omega}
\leq C_2h^m\|p\|_{m,\Omega}.
\]
类似地, 若 $\mathbf u\in H^{m+1}(\Omega)^N\cap V$, Hypothesis H1 给出
\[
\inf_{\mathbf v_h\in X_h}|\mathbf u-\mathbf v_h|_{1,\Omega}
\leq|\mathbf u-r_h\mathbf u|_{1,\Omega}
\leq C_3h^m\|\mathbf u\|_{m+1,\Omega}.
\]
这些不等式与 \eqref{eq:primitive-stokes-best-approximation} 蕴含
\eqref{eq:primitive-stokes-h1-pressure-error}.

还需证明极限 \eqref{eq:primitive-stokes-convergence}. 为此使用
Corollary~\ref{cor:primitive-general-mixed-convergence} 以及上述论证.
已知 $H^2(\Omega)\cap H_0^1(\Omega)$ 在 $H_0^1(\Omega)$ 中稠密,
且 $H^1(\Omega)\cap L_0^2(\Omega)$ 在 $L_0^2(\Omega)$ 中稠密.
因此, 使用上述算子 $r_h$ 和 $S_h$ 即可得到
\eqref{eq:primitive-stokes-convergence}.
\end{Proof}

当然, \eqref{eq:primitive-stokes-h1-pressure-error} 蕴含
$\|\mathbf u-\mathbf u_h\|_{0,\Omega}=O(h^m)$, 但可以使用
Theorem~\ref{thm:primitive-duality-error} 改进这一估计. 在这里取
\[
H=L^2(\Omega)^N.
\]
此时问题 \eqref{eq:primitive-dual-problem} 是齐次 Stokes 问题:
\begin{equation}\label{eq:primitive-stokes-dual-regularity-problem}
\left\{
\begin{aligned}
&\text{求 }(\boldsymbol\varphi,\xi)
\in H_0^1(\Omega)^N\times L_0^2(\Omega),\text{ 使得}\\
-\nu\Delta\boldsymbol\varphi+\operatorname{grad}\xi&=\mathbf g
&&\text{in }\Omega,\\
\operatorname{div}\boldsymbol\varphi&=0
&&\text{in }\Omega,
\end{aligned}
\right.
\qquad \mathbf g\in L^2(\Omega)^N.
\end{equation}
下面将需要该问题的正则性概念.

\setcounter{statement}{0}
\begin{Definition}\label{def:primitive-stokes-regularity}
若映射
\[
(\boldsymbol\varphi,\xi)
\longmapsto-\nu\Delta\boldsymbol\varphi+\operatorname{grad}\xi
\]
是从
\[
[H^2(\Omega)^N\cap V]
\times[H^1(\Omega)\cap L_0^2(\Omega)]
\]
到 $L^2(\Omega)^N$ 上的同构, 则称问题
\eqref{eq:primitive-stokes-dual-regularity-problem} 是正则的.
\end{Definition}

这一定义意味着, 每当右端 $\mathbf g\in L^2(\Omega)^N$ 时,
$\boldsymbol\varphi\in H^2(\Omega)^N$, $\xi\in H^1(\Omega)$, 且
\begin{equation}\label{eq:primitive-stokes-regularity-estimate}
\|\boldsymbol\varphi\|_{2,\Omega}+\|\xi\|_{1,\Omega}
\leq C\|\mathbf g\|_{0,\Omega}.
\end{equation}

\MathBookSourcePage{141}
由 Chapter~\ref{chap:stokes-foundations} 中 Stokes 正则性结果可知,
只要 $\Omega$ 的边界 $\Gamma$ 属于 $\mathcal C^2$ 类, 问题
\eqref{eq:primitive-stokes-dual-regularity-problem} 就是正则的.
当 $\Gamma$ 仅为 Lipschitz 连续时---后文中 $\Gamma$ 将是多边形线---,
若 $\Omega$ 是平面有界凸多边形, 同一问题仍然正则.

\setcounter{statement}{8}
\begin{Theorem}\label{thm:primitive-stokes-l2-error}
假设 Hypotheses H1, H2 和 H3 成立, 且问题
\eqref{eq:primitive-stokes-dual-regularity-problem} 正则. 若 Stokes 问题
\eqref{eq:primitive-homogeneous-stokes-strong} 的解 $(\mathbf u,p)$ 对某个满足
$1\leq m\leq l$ 的整数 $m$ 属于
\[
H^{m+1}(\Omega)^N\times(H^m(\Omega)\cap L_0^2(\Omega)),
\]
则有误差界
\begin{equation}\label{eq:primitive-stokes-l2-error}
\|\mathbf u-\mathbf u_h\|_{0,\Omega}
\leq Ch^{m+1}\left(
\|\mathbf u\|_{m+1,\Omega}+\|p\|_{m,\Omega}
\right).
\end{equation}
\end{Theorem}

\begin{Proof}
由 Theorem~\ref{thm:primitive-duality-error} 和
Remark~\ref{rem:primitive-duality-error-mixed-solution},
\begin{equation}\label{eq:primitive-stokes-duality-bound}
\begin{split}
\|\mathbf u-\mathbf u_h\|_{0,\Omega}
\leq{}&C_1\left\{
|\mathbf u-\mathbf u_h|_{1,\Omega}+\|p-p_h\|_{0,\Omega}
\right\}\\
&\times\sup_{\mathbf g\in L^2(\Omega)^N}
\frac{1}{\|\mathbf g\|_{0,\Omega}}
\left\{
\inf_{\boldsymbol\varphi_h\in X_h}
|\boldsymbol\varphi-\boldsymbol\varphi_h|_{1,\Omega}
+\inf_{\xi_h\in M_h}\|\xi-\xi_h\|_{0,\Omega}
\right\}.
\end{split}
\end{equation}
由于问题 \eqref{eq:primitive-stokes-dual-regularity-problem} 正则,
$\boldsymbol\varphi\in H^2(\Omega)^N\cap V$, 且
$\xi\in H^1(\Omega)\cap L_0^2(\Omega)$.
因此, Hypothesis H1 取 $m=1$, Hypothesis H2 取 $m=1$, 得到
\[
\inf_{\boldsymbol\varphi_h\in X_h}
|\boldsymbol\varphi-\boldsymbol\varphi_h|_{1,\Omega}
\leq C_3h|\boldsymbol\varphi|_{2,\Omega},
\]
\[
\inf_{\xi_h\in M_h}\|\xi-\xi_h\|_{0,\Omega}
\leq C_4h|\xi|_{1,\Omega}.
\]
将 \eqref{eq:primitive-stokes-regularity-estimate} 与
\eqref{eq:primitive-stokes-duality-bound} 结合, 得到
\[
\|\mathbf u-\mathbf u_h\|_{0,\Omega}
\leq C_5h\left\{
|\mathbf u-\mathbf u_h|_{1,\Omega}
+\inf_{q_h\in M_h}\|p-q_h\|_{0,\Omega}
\right\}.
\]
于是 \eqref{eq:primitive-stokes-l2-error} 由
\eqref{eq:primitive-stokes-h1-pressure-error} 和
\eqref{eq:primitive-h2-approximation} 得到.
\end{Proof}

如前一小节所述, 验证 Hypothesis H3 往往相当复杂.
事实上, 对空间 $X_h$ 和 $M_h$ 的选择通常取决于 H3 是否成立
(无论是猜想成立还是已经证明成立). 下一小节将说明如何构造这样的空间对.

最后简要回顾在 Subsection~\ref{subsec:primitive-decoupled-computation}
中提出的把 $\mathbf u_h$ 与 $p_h$ 的计算解耦的迭代方法.
取 $L^2(\Omega)$ 的标量积作为双线性形式 $c(\cdot,\cdot)$:
\[
c(p,q)=\int_\Omega p(x)q(x)\,dx.
\]
于是, 问题 \eqref{eq:primitive-stokes-discrete-kernel-problem} 的罚形式为:

求 $\mathbf u_h^\varepsilon\in X_h$, 使得
\begin{equation}\label{eq:primitive-stokes-penalty-problem}
\begin{split}
a(\mathbf u_h^\varepsilon,\mathbf v_h)
&+\frac{1}{\varepsilon}
(\rho_h(\operatorname{div}\mathbf u_h^\varepsilon),
\rho_h(\operatorname{div}\mathbf v_h))\\
&=\langle\mathbf f,\mathbf v_h\rangle
\qquad\forall\mathbf v_h\in X_h.
\end{split}
\end{equation}

\MathBookSourcePage{142}
这里 $\rho_h$ 是 $L^2(\Omega)$ 到 $Q_h$ 上的正交投影算子.
显然, 问题 \eqref{eq:primitive-stokes-penalty-problem} 把
$\mathbf u_h^\varepsilon$ 的计算与 $p_h^\varepsilon$ 的计算分离开来, 因为
\[
p_h^\varepsilon
=-\frac{1}{\varepsilon}\rho_h(\operatorname{div}\mathbf u_h^\varepsilon).
\]
当然, 只有当 $\rho_h(\operatorname{div}\mathbf v_h)$ 容易计算时,
该问题才有实际意义. 本章所讨论的几乎所有方法正是这种情形,
因为 $M_h$ 中的函数将是分片不连续的, 而 $\rho_h$ 是局部算子.

关于 $\mathbf u_h^\varepsilon$ 的收敛性, 直接应用
Theorem~\ref{thm:primitive-discrete-penalty-error} 得到下列结果.

\setcounter{statement}{9}
\begin{Theorem}\label{thm:primitive-stokes-penalty-error}
对所有 $\varepsilon>0$, 问题
\eqref{eq:primitive-stokes-penalty-problem} 有唯一解 $\mathbf u_h^\varepsilon$.
此外, 在 Hypothesis H3 下, 当 $\varepsilon\leq\varepsilon_0$ 足够小时,
\[
|\mathbf u_h-\mathbf u_h^\varepsilon|_{1,\Omega}
+\left\|p_h+\frac{1}{\varepsilon}
\rho_h(\operatorname{div}\mathbf u_h^\varepsilon)\right\|_{0,\Omega}
\leq C\varepsilon\|\mathbf f\|_{-1,\Omega},
\]
其中常数 $C>0$ 与 $h$ 和 $\varepsilon$ 无关.
\end{Theorem}

类似地, $\mathbf u_h^\varepsilon$ 和 $p_h^\varepsilon$ 可以按
$\varepsilon$ 的幂展开. 从 $p_h^0=p_h$ 开始, 定义解序列
$(\mathbf u_h^n,p_h^n)\in X_h\times M_h$:
\begin{equation}\label{eq:primitive-stokes-penalty-corrections}
\left\{
\begin{aligned}
a(\mathbf u_h^n,\mathbf v_h)
-(\operatorname{div}\mathbf v_h,p_h^n)&=0
&&\forall\mathbf v_h\in X_h,\\
(\operatorname{div}\mathbf u_h^n,q_h)
&=-(p_h^{n-1},q_h)
&&\forall q_h\in Q_h.
\end{aligned}
\right.
\end{equation}
Theorem~\ref{thm:primitive-discrete-penalty-expansion} 于是给出下列渐近展开.

\setcounter{statement}{10}
\begin{Theorem}\label{thm:primitive-stokes-penalty-expansion}
在 Hypothesis H3 下, 对所有整数 $M\geq1$, 当
$\varepsilon\leq\varepsilon_0$ 足够小时,
\[
\begin{split}
&\left|\mathbf u_h^\varepsilon-\mathbf u_h-
\sum_{n=1}^M\varepsilon^n\mathbf u_h^n\right|_{1,\Omega}\\
&\quad+\left\|\frac{1}{\varepsilon}
\rho_h(\operatorname{div}\mathbf u_h^\varepsilon)+p_h+
\sum_{n=1}^M\varepsilon^np_h^n\right\|_{0,\Omega}
\leq K_M\varepsilon^{M+1}\|\mathbf f\|_{-1,\Omega},
\end{split}
\]
其中常数 $K_M$ 与 $h$ 和 $\varepsilon$ 无关.
\end{Theorem}

现在讨论梯度算法. 在上述 $c(\cdot,\cdot)$ 的选择下,
双线性形式 $a_r^h(\cdot,\cdot)$ 为
\[
a_r^h(\mathbf u,\mathbf v)
=a(\mathbf u,\mathbf v)
+r(\rho_h(\operatorname{div}\mathbf u),
\rho_h(\operatorname{div}\mathbf v)).
\]
它在 $H_0^1(\Omega)^N$ 上椭圆, 且显然对称. 因此,
\eqref{eq:primitive-discrete-simple-gradient} 和
\eqref{eq:primitive-discrete-conjugate-gradient} 所描述的算法是真正的梯度算法.
带最优参数的简单梯度算法如下:

$1^\circ)$ 预测初值 $p_h^0\in M_h$, 并求 $\mathbf u_h^0\in X_h$:
\[
a_r^h(\mathbf u_h^0,\mathbf v_h)
=(p_h^0,\operatorname{div}\mathbf v_h)
+\langle\mathbf f,\mathbf v_h\rangle
\qquad\forall\mathbf v_h\in X_h.
\]

\MathBookSourcePage{143}
$2^\circ)$ 对 $m\geq0$, 已知 $(\mathbf u_h^m,p_h^m)$, 按下式确定
$\mathbf z_h^m\in X_h$, $\mu_h^m\in\mathbb R$ 以及
$(\mathbf u_h^{m+1},p_h^{m+1})\in X_h\times M_h$:
\[
a_r^h(\mathbf z_h^m,\mathbf v_h)
=-(\rho_h(\operatorname{div}\mathbf u_h^m),
\rho_h(\operatorname{div}\mathbf v_h))
\qquad\forall\mathbf v_h\in X_h,
\]
\[
\mu_h^m=-\frac{\|\rho_h(\operatorname{div}\mathbf u_h^m)\|_{0,\Omega}^2}
{(\rho_h(\operatorname{div}\mathbf u_h^m),
\rho_h(\operatorname{div}\mathbf z_h^m))},
\]
\[
p_h^{m+1}=p_h^m-\mu_h^m\rho_h(\operatorname{div}\mathbf u_h^m),
\qquad
\mathbf u_h^{m+1}=\mathbf u_h^m+\mu_h^m\mathbf z_h^m.
\]

共轭梯度算法按上式初始化 $\mathbf u_h^0$, 并以下列步骤代替第 $2^\circ)$ 步:

$2^\circ)$ 对 $m\geq0$, 已知 $(\mathbf u_h^m,p_h^m)\in X_h\times M_h$,
按下式计算 $(\mathbf z_h^m,\boldsymbol\omega_h^m)\in X_h\times M_h$,
$(\mu_h^m,\sigma_h^m)\in\mathbb R\times\mathbb R$ 以及
$(\mathbf u_h^{m+1},p_h^{m+1})\in X_h\times M_h$:
\[
\sigma_h^m=
\frac{\|\rho_h(\operatorname{div}\mathbf u_h^m)\|_{0,\Omega}^2}
{\|\rho_h(\operatorname{div}\mathbf u_h^{m-1})\|_{0,\Omega}^2}
\quad\text{only if $m\geq1$},
\]
\[
\boldsymbol\omega_h^m
=\rho_h(\operatorname{div}\mathbf u_h^m)
+\sigma_h^m\boldsymbol\omega_h^{m-1},
\qquad
\boldsymbol\omega_h^0=\rho_h(\operatorname{div}\mathbf u_h^0),
\]
\[
a_r^h(\mathbf z_h^m,\mathbf v_h)
=-(\boldsymbol\omega_h^m,\operatorname{div}\mathbf v_h)
\qquad\forall\mathbf v_h\in X_h,
\]
\[
\mu_h^m=-\frac{\|\rho_h(\operatorname{div}\mathbf u_h^m)\|_{0,\Omega}^2}
{(\rho_h(\operatorname{div}\mathbf u_h^m),
\rho_h(\operatorname{div}\mathbf z_h^m))},
\]
\[
p_h^{m+1}=p_h^m-\mu_h^m\boldsymbol\omega_h^m,
\qquad
\mathbf u_h^{m+1}=\mathbf u_h^m+\mu_h^m\mathbf z_h^m.
\]

由 Theorems~\ref{thm:primitive-discrete-simple-gradient-convergence}
和 \ref{thm:primitive-discrete-conjugate-gradient-convergence},
若 Hypothesis H3 成立, 两种梯度算法都收敛. 此外, 只要参数 $\mu_h^m$ 满足
\[
0<\inf_m\mu_h^m\leq\sup_m\mu_h^m
<2\left(\frac{\nu}{N}+r\right),
\]
简单梯度算法就收敛.

\subsection{inf-sup 条件的检验}\label{subsec:primitive-discrete-inf-sup-methods}

本小节构造一致满足 inf-sup 条件
\eqref{eq:primitive-discrete-inf-sup} 的空间对 $(X_h,M_h)$.
Boland \& Nicolaides [11] 提出的基本思想是: 若某个空间对
$(\bar X_h,\bar M_h)$ 一致满足该条件, 则只要另一些空间对满足局部 inf-sup 条件,
就可以生成一整族同样一致满足该条件的空间对. 换言之, 全局条件可以化为
更容易检验的局部条件.

\MathBookSourcePage{144}
设 $\Omega$ 被分成有限个互不相交的 Lipschitz 连续开子集
$\Omega_r$, 其边界为 $\Gamma_r$:
\[
\bar\Omega=\bigcup_{r=1}^R\bar\Omega_r.
\]
令 $X_h$ 和 $M_h$ 由 \eqref{eq:primitive-stokes-discrete-spaces} 定义,
且 $\mathbb R\subset Q_h$. 对 $1\leq r\leq R$, 置
\begin{equation}\label{eq:primitive-local-spaces}
\left\{
\begin{aligned}
X_h(\Omega_r)&=\{\mathbf v\in X_h;\ \mathbf v=\mathbf0
\text{ in }\Omega-\Omega_r\},\\
Q_h(\Omega_r)&=\{q|_{\Omega_r};\ q\in Q_h\},\\
M_h(\Omega_r)&=Q_h(\Omega_r)\cap L_0^2(\Omega_r),
\end{aligned}
\right.
\end{equation}
以及
\begin{equation}\label{eq:primitive-piecewise-constant-pressure-space}
\bar M_h=\{q\in L_0^2(\Omega);\ q|_{\Omega_r}\text{ 为常数},
\ 1\leq r\leq R\}.
\end{equation}
注意, $X_h(\Omega_r)$ 中的函数属于 $H_0^1(\Omega_r)^N$.
对于这一分割, 引入下列一致局部 inf-sup 条件作为假设.

\noindent\textbf{Hypothesis H4.}\quad
存在与 $h$ 和 $r$ 无关的常数 $\lambda^*>0$, 使得
\begin{equation}\label{eq:primitive-local-inf-sup}
\sup_{\mathbf v_h\in X_h(\Omega_r)}
\frac{\displaystyle\int_{\Omega_r}q_h\operatorname{div}\mathbf v_h\,dx}
{|\mathbf v_h|_{1,\Omega_r}}
\geq\lambda^*\|q_h\|_{0,\Omega_r}
\quad\forall q_h\in M_h(\Omega_r),\quad1\leq r\leq R.
\end{equation}

下面建立本小节的主要结果.

\setcounter{statement}{11}
\begin{Theorem}\label{thm:primitive-local-to-global-inf-sup}
设由 \eqref{eq:primitive-stokes-discrete-spaces} 定义的空间对 $(X_h,M_h)$
满足 Hypothesis H4. 若存在 $X_h$ 的子空间 $\bar X_h$, 使得空间对
$(\bar X_h,\bar M_h)$ 以与 $h$ 无关的常数 $\bar\beta$ 满足 inf-sup 条件
\eqref{eq:primitive-discrete-inf-sup}, 则 $(X_h,M_h)$ 也以与 $h$ 无关的常数
$\beta^*$ 满足该条件.
\end{Theorem}

\begin{Proof}
由 \eqref{eq:primitive-local-spaces} 立即得到 $Q_h(\Omega_r)$ 的正交分解
\[
Q_h(\Omega_r)=M_h(\Omega_r)\oplus\mathbb R.
\]
因此, 每个 $q_h\in M_h$ 可以分解为
\[
q_h=\tilde q_h+\bar q_h,
\]
其中
\[
\bar q_h|_{\Omega_r}
=\frac{1}{\operatorname{meas}(\Omega_r)}
\int_{\Omega_r}q_h\,dx,
\]
且 $\tilde q_r=\tilde q_h|_{\Omega_r}\in M_h(\Omega_r)$.
注意 $\bar q_h\in\bar M_h$, 而分解的正交性给出
\begin{equation}\label{eq:primitive-pressure-orthogonal-decomposition}
\|q_h\|_{0,\Omega}^2
=\|\tilde q_h\|_{0,\Omega}^2
+\|\bar q_h\|_{0,\Omega}^2.
\end{equation}

由 Hypothesis H4 和 Remark~\ref{rem:primitive-inf-sup-riesz-form},
存在 $\tilde{\mathbf v}_r\in X_h(\Omega_r)$, 使得
\begin{equation}\label{eq:primitive-local-divergence-lifting}
\left\{
\begin{aligned}
\int_{\Omega_r}\tilde q_r\operatorname{div}\tilde{\mathbf v}_r\,dx
&=\|\tilde q_r\|_{0,\Omega_r}^2,\\
|\tilde{\mathbf v}_r|_{1,\Omega_r}
&\leq\frac{1}{\lambda^*}\|\tilde q_r\|_{0,\Omega_r}.
\end{aligned}
\right.
\end{equation}

\MathBookSourcePage{145}
类似地, 因为空间对 $(\bar X_h,\bar M_h)$ 满足
\eqref{eq:primitive-discrete-inf-sup}, 存在 $\bar{\mathbf v}_h\in\bar X_h$, 使得
\begin{equation}\label{eq:primitive-coarse-divergence-lifting}
\left\{
\begin{aligned}
\int_\Omega\bar q_h\operatorname{div}\bar{\mathbf v}_h\,dx
&=\|\bar q_h\|_{0,\Omega}^2,\\
|\bar{\mathbf v}_h|_{1,\Omega}
&\leq\frac{1}{\bar\beta}\|\bar q_h\|_{0,\Omega}.
\end{aligned}
\right.
\end{equation}
令 $\tilde{\mathbf v}_h\in X_h$ 由
\[
\tilde{\mathbf v}_h|_{\Omega_r}=\tilde{\mathbf v}_r
\]
定义. 对某个 $\alpha>0$, 将函数
\[
\mathbf v_h=\tilde{\mathbf v}_h+\alpha\bar{\mathbf v}_h
\]
与 $q_h$ 相联系, 并调节参数 $\alpha$, 使 $(\mathbf v_h,q_h)$ 满足 inf-sup 条件.

估计 $(q_h,\operatorname{div}\mathbf v_h)$. 有
\[
\begin{split}
(q_h,\operatorname{div}\mathbf v_h)
={}&(\tilde q_h,\operatorname{div}\tilde{\mathbf v}_h)
+(\bar q_h,\operatorname{div}\tilde{\mathbf v}_h)\\
&+\alpha(\tilde q_h,\operatorname{div}\bar{\mathbf v}_h)
+\alpha(\bar q_h,\operatorname{div}\bar{\mathbf v}_h).
\end{split}
\]
这里,
\[
(\bar q_h,\operatorname{div}\tilde{\mathbf v}_h)=0,
\]
因为 $\tilde{\mathbf v}_r$ 在 $\Gamma_r$ 上为零; 并且由
\eqref{eq:primitive-local-divergence-lifting} 和
\eqref{eq:primitive-coarse-divergence-lifting}, 分别有
\[
(\tilde q_h,\operatorname{div}\tilde{\mathbf v}_h)
=\|\tilde q_h\|_{0,\Omega}^2,
\qquad
(\bar q_h,\operatorname{div}\bar{\mathbf v}_h)
=\|\bar q_h\|_{0,\Omega}^2.
\]
此外, 由 \eqref{eq:primitive-coarse-divergence-lifting},
\[
(\tilde q_h,\operatorname{div}\bar{\mathbf v}_h)
\leq\frac{\sqrt N}{\bar\beta}
\|\tilde q_h\|_{0,\Omega}\|\bar q_h\|_{0,\Omega}.
\]
汇总这些结果, 得到
\[
\begin{split}
(q_h,\operatorname{div}\mathbf v_h)
\geq{}&\|\tilde q_h\|_{0,\Omega}^2
+\alpha\|\bar q_h\|_{0,\Omega}^2\\
&-\alpha\frac{\sqrt N}{\bar\beta}
\|\tilde q_h\|_{0,\Omega}\|\bar q_h\|_{0,\Omega}.
\end{split}
\]
再用不等式
\[
\|\tilde q_h\|_{0,\Omega}\|\bar q_h\|_{0,\Omega}
\leq\varepsilon\|\tilde q_h\|_{0,\Omega}^2
+\frac{1}{4\varepsilon}\|\bar q_h\|_{0,\Omega}^2,
\]
得到
\[
\begin{split}
(q_h,\operatorname{div}\mathbf v_h)
\geq{}&\left(1-\frac{\alpha\varepsilon\sqrt N}{\bar\beta}\right)
\|\tilde q_h\|_{0,\Omega}^2\\
&+\alpha\left(1-\frac{\sqrt N}{4\varepsilon\bar\beta}\right)
\|\bar q_h\|_{0,\Omega}^2.
\end{split}
\]

\MathBookSourcePage{146}
这对任意 $\varepsilon>0$ 成立. 取
\[
\varepsilon=\frac12\frac{\bar\beta}{\alpha\sqrt N},
\]
则
\[
(q_h,\operatorname{div}\mathbf v_h)
\geq\frac12\|\tilde q_h\|_{0,\Omega}^2
+\alpha\left(1-\frac{\alpha N}{2\bar\beta^2}\right)
\|\bar q_h\|_{0,\Omega}^2.
\]
例如取
\[
\alpha=\frac{\bar\beta^2}{N}.
\]
由 \eqref{eq:primitive-pressure-orthogonal-decomposition} 得
\begin{equation}\label{eq:primitive-local-to-global-lower-bound}
(q_h,\operatorname{div}\mathbf v_h)
\geq\min\left(\frac12,\frac\alpha2\right)
\|q_h\|_{0,\Omega}^2.
\end{equation}
最后,
\[
\begin{split}
|\mathbf v_h|_{1,\Omega}
&\leq|\tilde{\mathbf v}_h|_{1,\Omega}
+\alpha|\bar{\mathbf v}_h|_{1,\Omega}\\
&\leq\frac{1}{\lambda^*}\|\tilde q_h\|_{0,\Omega}
+\frac{\alpha}{\bar\beta}\|\bar q_h\|_{0,\Omega}\\
&\leq\left\{
\left(\frac{1}{\lambda^*}\right)^2
+\left(\frac\alpha{\bar\beta}\right)^2
\right\}^{1/2}\|q_h\|_{0,\Omega},
\end{split}
\]
这里使用了 \eqref{eq:primitive-pressure-orthogonal-decomposition},
\eqref{eq:primitive-local-divergence-lifting} 和
\eqref{eq:primitive-coarse-divergence-lifting}. 将该不等式与
\eqref{eq:primitive-local-to-global-lower-bound} 结合, 得到所需的 inf-sup 条件
\[
\frac{(q_h,\operatorname{div}\mathbf v_h)}{|\mathbf v_h|_{1,\Omega}}
\geq\beta^*\|q_h\|_{0,\Omega},
\]
其中常数 $\beta^*$ 只依赖于 $N$, $\bar\beta$ 和 $\lambda^*$.
\end{Proof}

读者还可以参阅 Stenberg [75] 的相关方法.
